% sec:fd-dissipation — Kreiss--Oliger Dissipation
\section{Kreiss--Oliger Dissipation}
\label{sec:fd-dissipation}

High-order finite difference stencils converge rapidly for smooth data,
but they are susceptible to a subtle instability: grid-scale oscillations
--- parasitic modes with wavelength $2\,\Delta x$ --- can grow during a
time evolution even when the underlying solution is smooth.  These
oscillations arise from several sources: mesh refinement boundaries where
coarse-grid data is interpolated onto a fine grid, gauge dynamics that
introduce sharp features, or simply the nonlinear coupling in the BSSN
equations that transfers energy from resolved scales to the grid scale.
Once established, grid-scale noise is invisible to convergence testing
(it doesn't converge away --- it's a fixed-wavelength artefact of the
grid) and can destabilise the entire evolution.

Kreiss--Oliger dissipation addresses this by adding a small, high-order
diffusion term to the evolution equations.  The term is designed to damp
the highest-frequency modes aggressively while leaving well-resolved modes
virtually untouched, and --- crucially --- without degrading the
convergence order of the scheme \cite{Gustafsson1995}.
\coderef{dissipation}{rs}

\paragraph{The dissipation operator.}

The key observation is that grid-scale noise lives near the Nyquist
frequency $k_{\text{max}} = \pi/\Delta x$, while the physical solution
occupies much lower frequencies.  A dissipation operator that selectively
targets the Nyquist neighbourhood must act like a high-pass filter: large
response at the grid scale, negligible response at resolved scales.

The simplest such operator is the undivided even difference
$D_+^r D_-^r$, where $D_+$ and $D_-$ are the forward and backward
difference operators ($D_+ f_i = f_{i+1} - f_i$,
$D_- f_i = f_i - f_{i-1}$), and $r$ is a positive integer.  For $r = 1$
this is the familiar three-point second difference
$f_{i+1} - 2f_i + f_{i-1}$; higher $r$ produces a wider, more selective
filter.  The Kreiss--Oliger dissipation term added to the right-hand side
of the evolution equation $\partial_t f = \ldots$ is
%
\begin{equation}
  \epsilon_{\text{diss}}
  = (-1)^{r+1}\,\frac{\sigma}{\Delta x}\,
    \frac{D_+^r\,D_-^r\,f_i}{2^{2r}} \,,
  \label{eq:ko-operator}
\end{equation}
%
where $\sigma > 0$ is a dimensionless dissipation strength parameter and
the factor $2^{2r}$ normalises the operator so that its maximum value
(at the Nyquist frequency) is $\sigma/\Delta x$.  The sign
$(-1)^{r+1}$ ensures that the term always removes amplitude from
every Fourier mode, regardless of $r$ --- this is the defining property
of a dissipative operator and will be verified below.
\physnote{The dissipation term modifies the evolution equations at the
discrete level only.  It has no continuum counterpart: as
$\Delta x \to 0$, the term vanishes faster than the truncation error,
so it does not change the PDE being solved.}

\paragraph{Stencil coefficients.}

For a fourth-order spatial scheme, the dissipation order must be at least
fifth to avoid degrading the convergence rate.  Setting $r = 3$ gives a
fifth-order dissipation operator with a $2r + 1 = 7$-point stencil.
Expanding $D_+^3 D_-^3$ by repeated application of the three-point
second difference produces the stencil weights
%
\begin{equation}
  D_+^3\,D_-^3\,f_i = f_{i-3} - 6\,f_{i-2} + 15\,f_{i-1} - 20\,f_i
    + 15\,f_{i+1} - 6\,f_{i+2} + f_{i+3} \,,
  \label{eq:ko-stencil}
\end{equation}
%
with coefficients that are the binomial coefficients $\binom{6}{j}$
with alternating signs --- the sixth row of Pascal's triangle.  The
symmetry of the stencil follows from the symmetry of $D_+ D_-$: each
application of the three-point second difference preserves the centred
structure.
\implnote{The codebase function \texttt{ko\_diss\_x} implements the
seven-point stencil with coefficients $[-1, 6, -15, 20, -15, 6, -1]$
divided by $64$, matching \cref{eq:ko-stencil} divided by $2^{2r} = 64$
(up to the overall sign from $(-1)^{r+1}$, which is absorbed into
the calling convention).}

Including the sign $(-1)^{r+1} = (-1)^4 = 1$ and the normalisation, the
full dissipation contribution at each grid point is
%
\begin{equation}
\begin{split}
  \epsilon_{\text{diss}}
  = \frac{\sigma}{64\,\Delta x}
    \bigl(&f_{i-3} - 6\,f_{i-2} + 15\,f_{i-1} - 20\,f_i \\
    &+ 15\,f_{i+1} - 6\,f_{i+2} + f_{i+3}\bigr) \,.
\end{split}
  \label{eq:ko-explicit}
\end{equation}
%
The large negative central weight ($-20$) means the operator measures
how much $f_i$ deviates from a smooth interpolant through its six
neighbours.  For a polynomial of degree five or less, the stencil
returns exactly zero --- confirming that the operator annihilates the
smooth component of the solution.

\paragraph{Spectral analysis.}

To see why the operator damps grid-scale noise without affecting
resolved modes, consider its action on a Fourier mode
$f_j = e^{i k j \,\Delta x}$.  The forward and backward differences have
eigenvalues
%
\begin{equation}
  D_+ \to e^{ik\Delta x} - 1 \,, \qquad
  D_- \to 1 - e^{-ik\Delta x} \,,
  \label{eq:forward-backward-eigenvalues}
\end{equation}
%
and their product is
%
\begin{equation}
  D_+ D_- \to (e^{ik\Delta x} - 1)(1 - e^{-ik\Delta x})
  = 2\cos(k\,\Delta x) - 2 = -4\sin^2\!\bigl(\tfrac{k\,\Delta x}{2}\bigr) \,.
  \label{eq:dpdm-eigenvalue}
\end{equation}
%
The product is always non-positive: $D_+ D_-$ is a discrete diffusion
operator.  The $r$-th power therefore has eigenvalue
$(-1)^r\,4^r \sin^{2r}(k\,\Delta x/2)$, and the full dissipation
operator (\cref{eq:ko-operator}) acts on a Fourier mode as
%
\begin{equation}
  \epsilon_{\text{diss}} \to
  -\frac{\sigma}{\Delta x}\,\sin^{2r}\!\bigl(\tfrac{k\,\Delta x}{2}\bigr)
  \;\cdot\; f_j \,.
  \label{eq:ko-transfer}
\end{equation}
%
The transfer function $-(\sigma/\Delta x)\,\sin^{2r}(k\,\Delta x/2)$
is negative for all $k > 0$: every mode is damped, never amplified.
\warnnote{The negative sign in \cref{eq:ko-transfer} arises from the
product $(-1)^{r+1} \cdot (-1)^r = (-1)^{2r+1} = -1$: the sign
prefactor in \cref{eq:ko-operator} and the sign of the $(D_+D_-)^r$
eigenvalue combine to give a universally negative transfer function.
Getting this sign wrong produces anti-dissipation --- the operator
amplifies grid-scale noise instead of damping it.}

The damping rate depends sharply on frequency.  At the Nyquist limit
$k\,\Delta x = \pi$, $\sin^{2r}(\pi/2) = 1$ and the damping rate is
$\sigma/\Delta x$ --- the maximum.  For a well-resolved mode with
$k\,\Delta x \ll 1$, $\sin(k\,\Delta x/2) \approx k\,\Delta x/2$ and
the damping rate is approximately
$\sigma\,(k\,\Delta x/2)^{2r}/\Delta x = \sigma\,k^{2r}\,
\Delta x^{2r-1}/2^{2r}$, which vanishes as $\Delta x \to 0$.
For $r = 3$, the ratio of damping rates between the Nyquist mode and a
mode resolved by $10$ points per wavelength ($k\,\Delta x = 2\pi/10$) is
%
\begin{equation}
  \frac{\sin^6(\pi/2)}{\sin^6(\pi/10)}
  = \frac{1}{\sin^6(18°)}
  \approx \frac{1}{(0.309)^6}
  \approx 1150 \,.
  \label{eq:ko-selectivity}
\end{equation}
%
The operator damps the Nyquist mode over a thousand times more
aggressively than a mode resolved by ten grid points --- a selectivity
that justifies calling it a ``scalpel'' rather than a ``sledgehammer.''

\paragraph{Order preservation.}

Adding dissipation is worthless if it degrades the accuracy we worked to
achieve.  The dissipation must not degrade the $\order{\Delta x^q}$
convergence of the underlying spatial scheme.  Since the dissipation term
(\cref{eq:ko-operator}) scales as $\Delta x^{2r-1}$ for smooth modes
(the $\sin^{2r}$ factor contributes $(k\,\Delta x)^{2r}$ for fixed $k$),
the dissipation error is $\order{\Delta x^{2r-1}}$.  This is smaller
than the scheme's truncation error $\order{\Delta x^q}$ provided
$2r - 1 \geq q$, i.e.\ $r \geq (q + 1)/2$.  For a fourth-order scheme ($q = 4$),
$r \geq 3$ suffices, making the fifth-order ($r = 3$) dissipation operator
the minimum choice that preserves convergence order.
\xrefnote{The convergence tests of \cref{sec:fd-convergence} verify
that the dissipation term does not alter the measured convergence rate
at any tested resolution.}

\paragraph{Choice of $\sigma$.}

The dissipation strength $\sigma$ balances stability against accuracy.
Too small, and grid-scale noise survives long enough to contaminate the
evolution; too large, and the dissipation smears physical features.  For
the BSSN equations evolved with a fourth-order spatial scheme and a
Runge--Kutta time integrator, values of $\sigma$ in the range
$0.01$--$0.1$ are typical.  The codebase uses $\sigma = 0.05$ as the
default, which provides effective noise suppression across a wide range of
test problems without visibly affecting convergence.
\implnote{The function \texttt{apply\_ko\_dissipation} adds the
dissipation contribution in each coordinate direction to the right-hand
side for all $24$ evolved BSSN fields.  The stencil half-width of $3$
determines the minimum ghost-zone depth, as noted in
\cref{sec:fd-stencils}.}

The three-dimensional dissipation is simply the sum of the
one-dimensional operators applied in each coordinate direction.  No
cross-derivative terms appear: the operator is separable by construction.
On a uniform grid with $\Delta x = \Delta y = \Delta z$, the total
dissipation cost per grid point is three seven-point stencil
evaluations --- modest compared to the twelve stencil evaluations
required for the spatial derivatives of the BSSN right-hand side
(\cref{sec:fd-orders}).
